\chapter{Further consequences and future geometric extensions}
\label{chap:final-kervaire}

The current project ends at nonzero permanent survival of the standard
$h_6^2$. This chapter records further consequences of the paper. The geometric
statements below are future, out-of-scope sketches, not Challenge2 fields or
current acceptance obligations. No additional geometry root or proof module is
introduced by retaining these bibliographic statements.

\section{Classical two-line consequence}

% SOURCE: MainPaper/main.tex:151-163, Corollary cor:2line and its cited inputs.
\begin{corollary}[Survivors on the Adams two-line]
  \label{cor:adams-two-line-survivors}
  \uses{def:adams-hj,thm:external-adams-one-line,thm:h6-square-permanent}
  \notready
  Subject to the explicitly cited Ext calculations and differential results,
  the nonzero classes on the Adams two-line that survive to $E_\infty$ are
  precisely
  \[
    h_0h_2,\quad h_0h_3,\quad h_2h_4,\quad
    h_1h_j\ (j\ge3),\quad h_j^2\ (0\le j\le6).
  \]
\end{corollary}

\begin{proof}
  The supplied Adams, Lin, May, Mahowald--Tangora, Isaksen--Wang--Xu,
  Mahowald, Barratt--Jones--Mahowald, and Hill--Hopkins--Ravenel results first
  enumerate the $E_3$ candidates and then remove $h_2h_5$, $h_3h_5$,
  $h_3h_6$, and every $h_j^2$ with $j\ge7$.  Their positive survival results
  give every listed class except $h_6^2$; apply
  Theorem~\ref{thm:h6-square-permanent} for that final class.  The dependency
  points from this corollary to the main theorem, never conversely.
\end{proof}

\section{Future geometric hypotheses and endpoints (outside current scope)}

% SOURCE: MainPaper/main.tex:78--123; PROJECT_BOUNDARY.md, excluded geometry.
\begin{definition}[Additional geometric Kervaire hypotheses for future work]
  \label{def:geometric-kervaire-input}
  \uses{def:sphere-adams-coherence,def:framed-kervaire-context,thm:external-browder-criterion,thm:external-low-kervaire-existence,thm:external-hhr-nonexistence}
  \notready
  For a main input $I$, a future geometric extension would separately require
  an actual framed-manifold context $G$ and the Browder, low-dimensional
  existence and HHR results in that context. Its Pontryagin--Thom and Browder
  comparisons must refer to the same fixed sphere and classical Adams
  sequence as $I$, with explicit coherence when needed. This $G$ is not a
  field of current Challenge2 and has no current Lean delivery obligation.
\end{definition}

% SOURCE: MainPaper/main.tex:102-104, Theorem thm:main; future/out of scope.
\begin{theorem}[Conditional Kervaire manifold in dimension $126$]
  \label{thm:kervaire-126-conditional}
  \uses{def:kervaire-main-input,def:geometric-kervaire-input,thm:h6-square-permanent,thm:external-browder-criterion,def:kip126-trust-boundary}
  \notready
  For every main input $I$ and additional geometric hypotheses $G$ as above,
  there exists a closed framed smooth $126$-manifold in $G$'s context with
  Kervaire invariant one. This future conclusion does not follow from the
  current Challenge2 package alone.
\end{theorem}

\begin{proof}
  Apply Theorem~\ref{thm:h6-square-permanent} to $I$, obtaining survival of
  $h_6^2$, whose stem is $2^7-2=126$. Apply the permanence-to-existence direction of
  the separately supplied Browder theorem in $G$, after identifying its
  permanent-cycle premise with the fixed sphere sequence used in the first
  step. This is a future proof plan, not a current Lean construction.
\end{proof}

% SOURCE: MainPaper/main.tex:108-110, unlabelled Corollary 1.2; future/out of scope.
\begin{corollary}[Conditional exact list of Kervaire dimensions]
  \label{cor:kervaire-dimensions-exact}
  \uses{def:kervaire-main-input,def:geometric-kervaire-input,thm:kervaire-126-conditional,thm:external-low-kervaire-existence,thm:external-hhr-nonexistence,thm:external-browder-criterion}
  \notready
  For every main input $I$, additional geometric hypotheses $G$ as above,
  and natural number $n$, a framed smooth manifold in $G$'s context with
  Kervaire invariant one exists in
  dimension $n$ if and only if
  \[
    n=2\ \vee\ n=6\ \vee\ n=14\ \vee\ n=30\ \vee\ n=62\ \vee\ n=126.
  \]
\end{corollary}

\begin{proof}
  The separately supplied low-dimensional results in $G$ and the conditional
  $126$ theorem give existence in every listed dimension. The Browder theorem
  restricts any other positive dimension to $2^{j+1}-2$; the low-index fields
  cover $1\le j\le5$, the new theorem covers $j=6$, and the HHR result excludes
  every $j\ge7$. These additional geometric results are not current inputs.
\end{proof}
